\section{los big bucks}

Si pinchamos una esfera por un disco de cualquier dimensión, obtenemos una bola homológica.

\begin{teo}
    Dado un encaje $h:D^k\to \s^n$, la homología reducida es $\tilde{H}_i(\s^n - h(D^k))=0$ para todo $i\in\N$
\end{teo}

\comm{obviamente $k\leq n$ pero no se si decirlo, esto usa invariancia del dominio, no? q todo encaje va a cumplir $k\leq n$.}

\begin{dem}
    Hagamos inducción en $k\in\N$. 
    
    Si $k=0$, tenemos que $\s^n - h(D^0)$ es la esfera pinchada, es homeomorfa a $\R^n$, tomando un homeomorfismo que mande entornos reducidos del punto $h(D^0)$ en entornos del infinito. Con esto tenemos el caso base.

    Supongamos que se cumple para $k-1$, y probemos que entonces vale para $k$.

    Podemos buscar dos abiertos $A$ y $B$ convenientes tales que $X := \s^n-h(D^k) = A\cap B$ y usar la $\textbf{sucesión de}$ $\textbf{Mayer-Vietoris}$:
    
    $$\dots\rightarrow H_i(X)\rightarrow H_i(A)\oplus H_i(B)\rightarrow H_i(A\cup B)\rightarrow H_{i-1}(X)\rightarrow\dots$$
    
    Viendo $D^k$ como $I^k$, podemos tomar el cubrimiento de $\s^n-h\left(I^k\right)$ dado por los dos abiertos $A:=\s^n - h\left(I^{k-1}\times[0,\frac{1}{2}]\right)$ y $B:=\s^n - h\left(I^{k-1}\times[\frac{1}{2},1]\right)$.

    \comm{dibujito aca}

    En este caso $A\cup B = \s^n-h\left(I^{k-1}\times\{\frac{1}{2}\}\right)$. Esto lo podemos ver como $\s^n-\tilde{h}\left(D^{k-1}\right)$ con el encaje $\tilde{h}:D^{k-1}\to\s^n$ dado por precomponer $h$ con el homeomorfismo $I^{k-1}\times\{\frac{1}{2}\}\simeq D^{k-1}$. Aplicando la hipótesis de inducción obtenemos que $\tilde{H}_i(A\cup B)=0$ para todo $i\in\N$.

    Sustituyendo en la sucesión de Mayer-Vietoris, para cada $i\in\N$ tenemos la sucesión exacta

    $$0\rightarrow H_i(X)\xrightarrow{\Psi_i} H_i(A)\oplus H_i(B)\rightarrow 0$$

    Que esta sucesión sea exacta implica que $\Ker(\Psi_i)=0$, $\im(\Psi_i) = H_i(A)\oplus H_i(B)$ y por lo tanto $\Psi_i$ es un isomorfismo.

    Recordamos que los morfismos $\Psi_i$ de la sucesión de Mayer-Vietoris provienen de los morfismos entre los grupos $i$-cadenas singulares $\psi_i :C_i(X)\rightarrow C_i(A)\oplus C_i(B)$ definidos como $\psi_i(x)=(x,-x)$. 
        
    Por esta razón, a menos del signo, sabemos que las componentes de $\Psi_i$ están inducidas por las inclusiones $C_i(X)\hookrightarrow C_i(A)$, $C_i(X)\hookrightarrow C_i(B)$.

    Ahora queremos calcular explicitamente $H_i(X)$. El resultado que buscamos es que sea trivial, hagamos un razonamiento por absurdo.

    Supongamos que existe $x$ un representante de una clase no trivial del $i$-ésimo grupo de homología de $X$. Es decir, $x$ es un elemento de $C_i(X)$ tal que no es borde de un elemento en $C_{i+1}(X)$, o sea, $x\not\in\im(\partial_{i+1})$. La clase de $x$, al pasar por el isomorfismo $\Psi_i$, cae en un elemento no trivial de $H_i(A)\oplus H_i(B)$. Como $\Psi_i$ está inducida por la inclusión a menos de signos, la afirmación anterior implica que el ciclo $x$ no es un borde en al menos uno de $C_i(A)$ o $C_i(B)$.

    Sin pérdida de generalidad, asumamos que $x$ no es borde de ningún elemento de $C_i(A)$. Esto nos permite repetir el argumento tomando $X':=A=\s^n - h\left(I^{k-1}\times[0,\frac{1}{2}]\right)$, y eligiendo nuevos conjuntos $A':=\s^n - h\left(I^{k-1}\times[0,\frac{1}{4}]\right)$ y $B':=\s^n - h\left(I^{k-1}\times[\frac{1}{4},\frac{1}{2}]\right)$ dados por partir a la mitad el intervalo $[0,\frac{1}{2}]$ que definía al conjunto $A$. 
    
    Continuamos iterando sobre este argumento y obtenemos una sucesión de intervalos encajados $(I_m)_{m\in\N}$, tales que $I_0=[0,1]$, $I_j\subset I_{j-1}$, $\diam(I_m)\xrightarrow{m}0$.

    Por cómo fuimos definiendo la sucesión, se cumple que para todo $m\in\N$ el mismo $x$ que habíamos considerado antes va a ser un elemento de $C_i(\s^n - h(I^{k-1}\times I_m))$ que no es borde de ningún elemento de $C_{i+1}(\s^n - h(I^{k-1}\times I_m))$. Es decir, $x$ es un representante de una clase de homología no trivial de $\s^n - h(I^{k-1}\times I_m)$ para todo $m\in\N$.

    Sin embargo, por la hipótesis de inducción sabemos que $\tilde{H_i}(\s^n - h(I^{k-1}))=0$ para todo $i\in\N$, y por lo tanto, si llamamos $p:=\cap_{m=0}^\infty I_m$, tenemos que $\tilde{H_i}(\s^n - h(I^{k-1}\times\{p\}))=0$ para todo $i\in\N$. Esto implica que $x$ es borde de un elemento de $C_{i+1}(\s^n - h(I^{k-1}))$ al que llamaremos $y$.

    Los elementos de $C_{i+1}(\s^n - h(I^{k-1}))$ son sumas formales finitas de $(i+1)$-símplices singulares, es decir, sumas formales de finitas funciones continuas $\sigma:\Delta^{i+1}\to\s^n - h(I^{k-1}\times\{p\})$. Cada una de estas funciones tiene imagen compacta, por lo tanto la unión de las imágenes de las finitas funciones será un conjunto compacto.

    Considerando en particular el elemento $y$, la unión de las imágenes de sus símplices singulares es cubierta por la familia de conjuntos abiertos $\{\s^n - h(I^{k-1}\times I_m) : m\in\N\}$. Por ser compacta, existe un subcubrimiento finito. Como los conjuntos de la familia están encajados, esto implica que la unión de las imágenes de $y$ está contenida en $\s^n - h(I^{k-1}\times I_M)$ para algún $M\in\N$.
    
    Entonces, lo que obtuvimos al final es que $y$ es un elemento de $C_{i+1}(\s^n - h(I^{k-1}\times I_M))$ tal que $x=\partial_{i+1} y$. Esto es absurdo, ya que nos habíamos asegurado que $x$ fuera un representante de una clase de homología no trivial de $\s^n - h(I^{k-1}\times I_m)$.

    El absurdo surgió de tomar un elemento en una clase no trivial de $H_i(\s^n-h(D))$, por lo tanto concluimos que $H_i(\s^n-h(D))=0$.

    \begin{teo}
    Dado un encaje $h:\s^k\to \s^n$ con $k<n$, para todo $i\in\N$ se cumple que la homología reducida $\tilde{H}_i(\s^n - h(\s^k))$ es $\Z$ si $i=n-k-1$ y $0$ en cualquier otro caso.
    \end{teo}

    \begin{dem}

    Para probar esto vamos a apoyarnos en la parte anterior. La idea es hacer la prueba por inducción en $k$ y usar la sucesión de Mayer-Vietoris con un cubrimiento por abiertos de la forma $\s^n - h(D^k)$. 

    Trivialmente tenemos el caso $k=0$. Como $\s^n - h(\s^0)$ es la esfera con dos pinchaduras, homeomorfa a $\s^{n-1}\times\R$, se cumple que $H_i(\s^n - h(\s^0))=0$ en todo caso salvo cuando $i=0$, y ahí vale $H_0(\s^n - h(\s^0))=\Z\oplus\Z$, y por lo tanto la homología reducida es $\tilde{H_0}(\s^n - h(\s^0))=\Z$.

    Asumamos ahora que vale para $k-1$. Sean $D^k_+$, $D^k_-$ dos discos cerrados de dimensión $k$ tales que $D^k_+\cap D^k_- \simeq\s^{k-1}$. Por la parte anterior sabemos que si definimos $A:=\s^k-h(D^k_-)$, $B:=\s^k-h(D^k_-)$, entonces se cumple que $\tilde{H_i}(A)=\tilde{H_i}(B)=0$ para todo $i\in\N$.

    Viendo la sucesión de Mayer-Vietoris para $\s^n- h(\s^k)$ cubierto por $A$ y $B$, el término $H_i(A)\oplus H_i(B)$ se anula para todo $i\in\N$, entonces obtenemos las siguientes sucesiones exactas:

    $$0\rightarrow H_i(\s^n- h(\s^k))\rightarrow H_{i-1}(A\cap B)\rightarrow 0$$



    \end{dem}


    
 

    % Un cubrimiento de $\s^n-h(I^k)$ está dado por $\s^n-A\cap\s^n-B$, entonces por 

    % $\s^n -h(I^k)= \s^n -h(A\cup B) = \s^n-h(A) \cap \s^n-h(B)$
    
\end{dem}